%% LyX 2.5.1 created this file.  For more info, see https://www.lyx.org/.
%% Do not edit unless you really know what you are doing.
\documentclass[12pt,twoside]{NCEPUDoctor}
\setcounter{secnumdepth}{3}
\setcounter{tocdepth}{3}
\usepackage[table]{xcolor}
\usepackage{colortbl}
\usepackage{booktabs}
\usepackage{amsmath}
\usepackage{amssymb}
\usepackage{graphicx}
\usepackage{fancyhdr}
\pagestyle{fancy}

\makeatletter

%%%%%%%%%%%%%%%%%%%%%%%%%%%%%% LyX specific LaTeX commands.
\DeclareTextSymbolDefault{\textquotedbl}{T1}
%% Because html converters don't know tabularnewline
\providecommand{\tabularnewline}{\\}

\@ifundefined{date}{}{\date{}}
%%%%%%%%%%%%%%%%%%%%%%%%%%%%%% User specified LaTeX commands.
% ============================================================
% 额外宏包
\usepackage{multirow}
\usepackage{makecell}
\usepackage{adjustbox}

% ===== 列表间距 =====
\usepackage{enumitem}
\renewcommand{\labelenumi}{[\arabic{enumi}]}
\setenumerate[1]{itemsep=0pt,partopsep=0pt,parsep=0pt,topsep=0pt,leftmargin=2.5em,labelwidth=2em,labelsep=0.5em}
\setitemize[1]{itemsep=0pt,partopsep=0pt,parsep=\parskip,topsep=5pt}
\setdescription{itemsep=0pt,partopsep=0pt,parsep=\parskip,topsep=5pt}

% ===== 首段缩进 =====
\usepackage{indentfirst}
\setlength{\parindent}{2em}

% ===== LyX 摘要辅助命令 =====
\newenvironment{lyxabstractzh}{%
  \zihao{-4}\sloppy
}{}
\newcommand{\lyxzhkeywords}[1]{%
  \par\vspace{1em}
  \noindent\textbf{关键词：} #1
  \clearpage
}
\newenvironment{lyxabstracten}{%
  \zihao{-4}\sloppy
  \hyphenpenalty=10000
}{}
\newcommand{\lyxenkeywords}[1]{%
  \par\vspace{1em}
  \noindent\textbf{Keywords:} #1
  \clearpage
}

% ===== 防止 LyX 自动生成 \maketitle =====
\renewcommand{\maketitle}{}

% ===== 算法 LyX 支持 =====
\newcommand{\algpseudobegin}{\begin{algorithmic}[1]}
\newcommand{\algpseudoend}{\end{algorithmic}}
\newcommand{\algstateinner}[1]{\STATE #1}
\newcommand{\algendforinner}[1]{\ENDFOR}
\newcommand{\algendifinner}[1]{\ENDIF}
\newcommand{\algendwhileinner}[1]{\ENDWHILE}
\newcommand{\algelseinner}[1]{\ELSE}
\newcommand{\algforinner}[1]{\FOR{#1}}
\newcommand{\algifinner}[1]{\IF{#1}}
\newcommand{\algwhileinner}[1]{\WHILE{#1}}

% ===== 参考文献标题间距修复 =====
\makeatletter
\renewcommand{\bibsection}{%
  \chapter*{\bibname}%
  \@mkboth{\bibname}{\bibname}%
  \addcontentsline{toc}{chapter}{\bibname}%
  \vspace{0.5\baselineskip}%
}
\makeatother

% ===== 算法环境 =====
\makeatletter
\renewcommand{\fps@algorithm}{!t}
\makeatother

% ===== 前言页码 =====
\makeatletter
\let\origcontents\contents
\renewcommand{\contents}{%
  \clearpage
  \tableofcontents
  \pagestyle{plain}%
}
\makeatother

% ===== 表格标题 =====
\captionsetup[table]{labelsep=space}

% ===== 引用上标 =====
\setcitestyle{super}

\makeatother

\begin{document}
\degreetype{academic}
\ctitlezh{智能电网关键关键技术研究与应用}
\ctitleen{Research on Power System Fault Diagnosis Based on Deep Learning}
\Ltitle{False}
\cauthorzh{张三}
\cauthoren{Zhang San}
\cadvisorzh{李教授}
\cadvisoren{Prof. Li}
\ccosupervisoren{Prof. Wang}
\cdegreezh{工学博士}
\cdegreeen{Doctor of Engineering}
\cschoolzh{电气与电子工程学院}
\cschoolen{School of Electrical and Electronic Engineering}
\cdatezh{2024年6月}
\cdateen{June, 2024}
\csubjectzh{电气工程}
\csubjecten{Electrical Engineering}
\cmajorzh{电力系统及其自动化}
\cmajoren{Power System and Its Automation}
\cclc{xxxx}
\cudc{xxxx}
\csecretlevel{public}

\maketitle
\maketitlepage
\makechinainfopage
\makeenglishinfopage
\makestatement
\pagenumbering{Roman}
\setcounter{page}{0}

\chapterstartdual{Abstract}{摘\qquad{}要}
\begin{lyxabstractzh}
摘要是论文内容的高度概括，应具有独立性和自含性，即不阅读论文的全文，就能获得必要的信息。摘要应包括本论文的目的、主要研究内容、研究方法、创造性成果及其理论与实际意义。摘要中不宜使用公式、化学结构式、图表和非公知公用的符号和术语，不标注引用文献编号。避免将摘要写成目录式的内容介绍。
\end{lyxabstractzh}

\lyxzhkeywords{关键词1；关键词2；关键词3；关键词4；关键词5}

\chapterstartdual{Abstract}{Abstract}
\begin{lyxabstracten}
The key contributions of this thesis include the development of a
new model for predicting the load capacity of gas bearings, the optimization
of bearing design parameters for specific applications, and the experimental
validation of the theoretical findings. The findings of this research
have significant implications for the design and application of gas
bearings in various industries.
\end{lyxabstracten}

\lyxenkeywords{keyword 1, keyword 2, keyword 3, keyword 4, keyword 5}

\contents{}

\listoffiguresandtables{}

\chapterstartdual{Nomenclature}{主要符号表}

\begin{longtable}{ C{0.4} C{0.6}}
\toprule[1.5pt]
符号 & 含义 \\
\midrule[1pt]
\endfirsthead
\toprule[1.5pt]
符号 & 含义 \\
\midrule[1pt]
\endhead
\bottomrule[1.5pt]
\endfoot
0.35 & 0.136 \\
0.40 & 0.171 \\
0.35 & 0.136 \\
0.21753 & 246900 \\
5.39252 & -4.66248 \\
0.40 & 0.171 \\
0.25485 & 296900 \\
5.47261 & -4.59372 \\
0.35 & 0.136 \\
0.21753 & 246900 \\
5.39252 & -4.66248 \\
0.40 & 0.171 \\
0.25485 & 296900 \\
5.47261 & -4.59372 \\
0.35 & 0.136 \\
0.40 & 0.171 \\
0.35 & 0.136 \\
0.21753 & 246900 \\
5.39252 & -4.66248 \\
0.40 & 0.171 \\
0.25485 & 296900 \\
\end{longtable}

\clearpage\setcounter{page}{1}\pagenumbering{arabic}

\chapterdual{Introduction}{绪论}

\sectiondual{Background, objective and significance of the subject}{课题背景及研究的目的和意义}

发展国防工业微电子工业等尖端技术，需要精密和超精密的仪器设备。精密仪器设备要求高速、…… （宋体、Times New Roman 小4号字）

\sectiondual{Developmental of gas-lubricated bearing and correlated theories}{气体润滑轴承及其相关理论的发展概况}

气体轴承是利用气膜支撑负荷或减少摩擦的机械构件。……

\subsectiondual{Developmental of gas-lubricated bearing}{气体润滑轴承的发展}

气体润滑轴承的发展1828年，R.R.Willis发表了一篇关于小孔节流平板中压力分布的文章，这是有记载的研究气体润滑的最早文献。……华北电力大学网信处

\subsectiondual{Research on static characteristics of porous externally Pressurized
gas bearing}{多孔质气体静压轴承的研究}

由于气体的压力低和可压缩性，……。

\begin{equation}
C^{f}_{t}=\left\{ a(P^{gen}_{t})^{2}+bP^{gen}_{t}+c\right\} \bigtriangleup t\label{gt1}
\end{equation}


\subsubsection{多孔质静压轴承的分类}


\subsubsection{多孔质材料特性的研究}

材料的主要特点是具有一定的……。

（1）孔隙特性 多孔质材料是由……。

\sectiondual{Computational fluid dynamics and applications of correlated field}{计算流体动力学及其在相关领域中的应用}


\chapterdual{Introduction}{绪论}

公式原则上应居中书写。若公式前有文字（如\textquotedbl 解\textquotedbl 、\textquotedbl 假定\textquotedbl 等），文字空两格写，公式仍居中写。公式末不加标点。公式序号按章编排，如第1章第一个公式序号为\textquotedbl (1-1)\textquotedbl 。文中引用公式时，一般用\textquotedbl 见式
(1-1)\textquotedbl 或\textquotedbl 由公式 (1-1)\textquotedbl 。公式中用斜线表示\textquotedbl 除\textquotedbl 的关系时应采用括号，以免含糊不清。通常\textquotedbl 乘\textquotedbl 的关系在前。不能用文字形式表示等式。

示例引用公式\ref{gt2}。

\begin{gather}
C^{f}_{t}=\left\{ a(P^{gen}_{t})^{2}+bP^{gen}_{t}+c\right\} \bigtriangleup t\label{gt2}\\
C^{f}_{t}=\left\{ a(P^{gen}_{t})^{2}+bP^{gen}_{t}+c\right\} \bigtriangleup t\label{gt3}
\end{gather}

式中\textasciitilde\textasciitilde  $P_t^{gen}$------ 发电量；\breakwithtilde $P_t^{gen}$------ 发电量；\breakwithtilde $P_t^{gen}$------ 发电量。

研究背景利用现成的商用软件来研究流场，可以免去对N-S方程求解程序的研究背景利用现成的商用软件来研究流场，

\begin{equation}
C^{f}_{t}=\left\{ a(P^{gen}_{t})^{2}+bP^{gen}_{t}+c\right\} \bigtriangleup t\label{gt1b}
\end{equation}


\chapterdual{Introduction}{绪论}

\sectiondual{示例章节}{示例章节}

本节演示本模板中各类数学环境与算法环境的使用方法。
\begin{definition}[向量范数]
设$\mathbf{x}=(x_{1},\dots,x_{n})\in\mathbb{R}^{n}$，向量的$2$范数定义为
\[
\|\mathbf{x}\|_{2}=\sqrt{\sum^{n}_{i=1}x^{2}_{i}}.
\]
\end{definition}
\begin{theorem}[Cauchy-{}-Schwarz 不等式]
对于任意$\mathbf{x},\mathbf{y}\in\mathbb{R}^{n}$，成立：
\[
|\mathbf{x}^{\top}\mathbf{y}|\le\|\mathbf{x}\|_{2}\,\|\mathbf{y}\|_{2}.
\]
\end{theorem}
\begin{lemma}
若向量$\mathbf{x}$与$\mathbf{y}$线性无关，则$\|\alpha\mathbf{x}+\beta\mathbf{y}\|_{2}=0$当且仅当$\alpha=\beta=0$。
\end{lemma}
\begin{corollary}
若$\mathbf{x}\neq\mathbf{0}$，则有$\|\mathbf{x}\|_{2}>0$。
\end{corollary}
\begin{proposition}
对任意标量$c$，向量范数满足齐次性：$\|c\mathbf{x}\|_{2}=|c|\|\mathbf{x}\|_{2}$。
\end{proposition}
\begin{example}
例如，向量$\mathbf{x}=(3,4)$的$2$范数为$\|\mathbf{x}\|_{2}=5$。
\end{example}
\begin{remark}
上述性质表明二范数是欧氏空间中最常用的度量方式。
\end{remark}
\begin{proof}
不等式的证明可由平方展开并利用非负性条件得到。具体地，
\[
0\le\|\mathbf{x}-t\mathbf{y}\|^{2}_{2}
\]
展开并对
\end{proof}
算法引用示例\algref{alg:population-sim}。

\begin{algorithm}[H]  \caption{基于个体模拟的种群动态演化算法}  \label{alg:population-sim}  \begin{algorithmic}[1]    \STATE 初始化种群 $\mathcal{P}_0$，包含 $N$ 个个体，每个个体具有特征向量 $\mathbf{x}_i \in \mathbb{R}^d$    \STATE 初始化环境承载容量 $K$，基础繁殖率 $\lambda$，变异幅度 $\sigma$，自然淘汰阈值 $\epsilon$    \STATE 初始化资源分布矩阵 $\mathbf{R} \in \mathbb{R}^{M \times M}$，季节周期 $T_s$    \FOR{$t = 1$ 至 $T_{max}$}      \STATE 根据季节因子更新资源分布：$\mathbf{R}_t \gets \mathbf{R} \cdot (1 + \beta \sin(2\pi t / T_s))$      \STATE 计算种群密度 $\rho_t = |\mathcal{P}_t| / K$      \WHILE{存在未评估的个体}        \STATE 计算个体 $i$ 的适应度 $f_i = g(\mathbf{x}_i, \mathbf{R}_t) \cdot (1 - \rho_t)$        \IF{$f_i < \epsilon$}          \STATE 将个体 $i$ 标记为淘汰，从种群中移除        \ELSE          \STATE 以概率 $p_i = \lambda \cdot f_i / \bar{f}$ 判定个体 $i$ 是否繁殖          \IF{个体 $i$ 进行繁殖}            \STATE 随机选择配偶 $j$，要求 $\|\mathbf{x}_i - \mathbf{x}_j\| \leq \delta$            \STATE 生成后代特征：$\mathbf{x}_{new} = \frac{1}{2}(\mathbf{x}_i + \mathbf{x}_j) + \mathcal{N}(0, \sigma^2 \mathbf{I})$            \STATE 将后代加入下一代种群 $\mathcal{P}_{t+1}$          \ENDIF        \ENDIF      \ENDWHILE      \FOR{每个存活个体 $i \in \mathcal{P}_{t+1}$}        \STATE 根据资源梯度更新位置：$\mathbf{x}_i \gets \mathbf{x}_i + \eta \, \nabla_{\mathbf{x}} g(\mathbf{x}_i, \mathbf{R}_t)$      \ENDFOR    \ENDFOR  \end{algorithmic}\end{algorithm}

\chapterdual{Experimental Study on Local Porous Hydrostatic Bearings}{基于FLUENT软件的轴承静态特性研究}

\sectiondual{Analysis of FLUENT simulation results}{FLUENT仿真结果分析}

图\ref{fig:system}表示出在局部多孔质圆柱塞的半径不同时，局部多孔质圆柱塞内和气膜内的压力分布变化情况。其中的半径分别是rr=1.5mm，2.5mm，3.5mm，4.5mm。孔质圆柱塞的半径不同时，局部多孔质圆柱塞内和气膜内的压力分布变化情况。

\begin{figure}
\includegraphics[width=0.8\linewidth]{figures/1}

\caption{局部多孔质圆柱塞半径不同时轴承的压力分布云图}\label{fig:system}
\end{figure}

从图\ref{fig:bearing-pressure}可知，由于节流半径不同，导致节流效果产生很大的不同，其中，半径小的节流效果明显，图\subref{subfig:bearing-15}对应的压力变化最明显，而图\subref{subfig:bearing-45}的变化非常小，这导致了气膜内的压力分布产生了很大的不同。从而使承载能力随着半径的增加而得到很大的提高。从图\ref{fig:bearing-pressure}可知，由于节流半径不同，导致节流效果产生很大的不同，其中，半径小的节流效果明显，图\subref{subfig:bearing-15}对应的压力变化最明显，而图\subref{subfig:bearing-45}的变化非常小，这导致了气膜内的压力分布产生了很大的不同。从而使承载能力随着半径的增加而得到很大的提高。从图\ref{fig:bearing-pressure}可知，由于节流半径不同，导致节流效果产生很大的不同，其中，半径小的节流效果明显，图\subref{subfig:bearing-15}对应的压力变化最明显，而图\subref{subfig:bearing-45}的变化非常小，这导致了气膜内的压力分布产生了很大的不同。从而使承载能力随着半径的增加而得到很大的提高。从图\ref{fig:bearing-pressure}可知，由于节流半径不同，导致节流效果产生很大的不同，其中，半径小的节流效果明显，图\subref{subfig:bearing-15}对应的压力变化最明显，而图\subref{subfig:bearing-45}的变化非常小，这导致了气膜内的压力分布产生了很大的不同。从而使承载能力随着半径的增加而得到很大的提高。

\begin{figure}
\begin{subfigure}[b]{0.48\textwidth}\centering\includegraphics[width=\textwidth]{figures/1.png}\subfigcaptioncneng{$R_3$=1.5mm 时轴承的压力分布云图。}                      {Pressure contour of bearing when $R_3$=1.5mm。}\label{subfig:bearing-15}\end{subfigure}\hfill\begin{subfigure}[b]{0.48\textwidth}\centering\includegraphics[width=\textwidth]{figures/1.png}\subfigcaptioncneng{$R_3$=2.5mm 时轴承的压力分布云图。}                      {Pressure contour of bearing when $R_3$=2.5mm。}\label{subfig:bearing-25}\end{subfigure}\vspace{1em}\begin{subfigure}[b]{0.48\textwidth}\centering\includegraphics[width=\textwidth]{figures/1.png}\subfigcaptioncneng{$R_3$=3.5mm 时轴承的压力分布云图。}                      {Pressure contour of bearing when $R_3$=3.5mm。}\label{subfig:bearing-35}\end{subfigure}\hfill\begin{subfigure}[b]{0.48\textwidth}\centering\includegraphics[width=\textwidth]{figures/1.png}\subfigcaptioncneng{$R_3$=4.5mm 时轴承的压力分布云图。}                      {Pressure contour of bearing when $R_3$=4.5mm。}\label{subfig:bearing-45}\end{subfigure}

\caption{\$R\_3\$=1.5mm 时轴承的压力分布云图。}\label{fig:bearing-pressure}
\end{figure}

从图\ref{fig:bearing-pressure}可知，由于节流半径不同，导致节流效果产生很大的不同，其中，半径小的节流效果明显，图\subref{subfig:bearing-15}对应的压力变化最明显，而图\subref{subfig:bearing-45}的变化非常小，这导致了气膜内的压力分布产生了很大的不同。从而使承载能力随着半径的增加而得到很大的提高。

\chapterdual{Experimental Study on Local Porous Hydrostatic Bearings}{局部多孔质静压轴承的实验研究}

学位论文的结论作为论文正文的最后一章单独排写\cite{clerc2010discrete}。

结论是对整个论文主要成果的总结。在结论中应明确指出本研究内容的创造性成果或创新点理论（含新见解、新观点），并指出今后进一步在本研究方向进行研究工作的展望与设想，不要将结论写成论文的摘要。结论内容一般在2000字以内。

\begin{table}
\caption{1号试样渗透率测试数据（温度： T=16℃ 高度： H=5.31mm）}\label{tab:notation}

\centering{}%
\begin{tabular}{C{0.4}C{0.6}}
\toprule 
Notation & Description\tabularnewline
$P^{L}_{t}$ & Users' energy consumption\tabularnewline
$P^{gen}_{t}$ & The power of generator\tabularnewline
\bottomrule
\end{tabular}
\end{table}

结论是对整个论文主要成果的总结。在结论中应明确指出本研究内容的创造性成果或创新点理论（含新见解、新观点），并指出今后进一步在本研究方向进行研究工作的展望与设想，不要将结论写成论文的摘要。结论内容一般在2000字以内。结论是对整个论文主要成果的总结。在结论中应明确指出本研究内容的创造性成果或创新点理论（含新见解、新观点），并指出今后进一步在本研究方向进行研究工作的展望与设想，不要将结论写成论文的摘要。结论内容一般在2000字以内。

\begin{longtable}{C{0.15} C{0.15} C{0.25} C{0.15} C{0.15} C{0.15}}\bilingualtablecaption{1号试样渗透率测试数据（温度：T=16  高度：H=5.31mm）}{Data of measured permeability of sample 1 (Temperature: T=16 Height: H=5.31mm)}\label{tab:permeability} \\\toprule[1.5pt]供气压力 & 流量测量 & 流量修正值 & 压力差 & $\lg \Delta P$ & $\lg M$ \\$P_g$ (MPa) & $M'$ (m$^3$/h) & $M$ ($\times10^{-4}$ m$^3$/s) & $\Delta P$ (Pa) & & \\\midrule[1pt]\endfirsthead\multicolumn{6}{r}{{\thenexttable}} \\\toprule[1.5pt]供气压力 & 流量测量 & 流量修正值 & 压力差 & $\lg \Delta P$ & $\lg M$ \\$P_g$ (MPa) & $M'$ (m$^3$/h) & $M$ ($\times10^{-4}$ m$^3$/s) & $\Delta P$ (Pa) & & \\\midrule[1pt]\endhead\bottomrule[1.5pt]\endfoot0.35 & 0.136 & 0.21753 & 246900 & 5.39252 & -4.66248 \\0.40 & 0.171 & 0.25485 & 296900 & 5.47261 & -4.59372 \\0.35 & 0.136 & 0.21753 & 246900 & 5.39252 & -4.66248 \\0.40 & 0.171 & 0.25485 & 296900 & 5.47261 & -4.59372 \\0.35 & 0.136 & 0.21753 & 246900 & 5.39252 & -4.66248 \\0.40 & 0.171 & 0.25485 & 296900 & 5.47261 & -4.59372 \\0.35 & 0.136 & 0.21753 & 246900 & 5.39252 & -4.66248 \\0.40 & 0.171 & 0.25485 & 296900 & 5.47261 & -4.59372 \\0.35 & 0.136 & 0.21753 & 246900 & 5.39252 & -4.66248 \\0.40 & 0.171 & 0.25485 & 296900 & 5.47261 & -4.59372 \\0.45 & 0.202 & 0.28467 & 346900 & 5.54020 & ——\end{longtable}

\chapterdual{Conclusions and prospects}{结论与展望}

学位论文的结论作为论文正文的最后一章单独排写\cite{chen2001hao}。\cite{clerc2010discrete,liu1979,maoxia,xia2002analysis,1025090578.nh}，\cite{xia2002analysis,1025090578.nh,xxx2023},\cite{award2007,gbt16159,maoxia2000,feng1997,wang1998}。

结论是对整个论文主要成果的总结。在结论中应明确指出本研究内容的创造性成果或创新点理论（含新见解、新观点），并指出今后进一步在本研究方向进行研究工作的展望与设想，不要将结论写成论文的摘要。结论内容一般在2000字以内。学位论文的结论作为论文正文的最后一章单独排

\clearpage

\chapterstartdual{Reference}{参考文献}

\bibliographystyle{tools/gbt7714-numerical}
\bibliography{reference}

\clearpage

\startappendix{}

\chapterdual{Example of Appendix Title}{附录标题示例}

附录作为主体部分的补充，并不是必须的。

下列内容可以作为附录编于论文后：

------为了整篇论文材料的完整，但编入正文又有损于编排的条理性和逻辑性，这一材料包括比正文更为详尽的信息、研究方法和技术更深入的叙述，对了解正文内容有用的补充信息等；

------由于篇幅过大或取材于复制品而不便于编入正文的材料；

------不便于编入正文的罕见珍贵资料；

------对一般读者并非必要阅读，但对本专业同行有参考价值的资料；

------某些重要的原始数据、数学推导、结构图、统计表、计算机打印输出件等。

\stopappendix{}

\clearpage

\chapterstartdual{Papers published in the period of Ph.D. education}{攻读博士学位期间发表的论文及其它成果}

\boldparafirst{（一）发表的学术论文}
\begin{enumerate}
\item xxx，xxx．局部多孔质气体静压轴向轴承静态特性的数值求解{[}J{]}．摩擦学学报，2007，38(12)：68-{}-72．（EI收录号：20071510544816）
\item xxx，xxx．基于深度学习的电力系统故障诊断方法研究{[}J{]}．中国电机工程学报，2023，43(5)：1024-{}-1035．
\end{enumerate}

\boldpara{（二）申请及已获得的专利}
\begin{enumerate}
\item xxx，xxx. 一种温热外敷药制备方案：中国，88105607.3{[}P{]}. 1989-07-26．
\end{enumerate}

\boldpara{（三）获得的科技奖励}
\begin{enumerate}
\item xxx，xxx. xx静载下预应力混凝土房屋结构设计统一理论. 黑龙江省科学技术二等奖, 2007.
\end{enumerate}
注：如已发表的学术论文被EI或SCI收录，请标明收录号及SCI论文的影响因子；对已接收但尚未发表出来的学术论文，请注明是否EI或SCI刊源。

\clearpage

\chapterstartdual{Research Work in the period of Ph.D. education}{攻读博士学位期间参加的科研工作}
\begin{enumerate}
\item 国家自然科学基金项目（项目编号：xxxxxxxx），电力系统故障诊断方法研究。
\item 国家重点研发计划项目（项目编号：xxxxxxxx），智能电网关键技术与装备研发。
\end{enumerate}
\clearpage

\chapterstartdual[致谢]{Acknowledgements}{致\textasciitilde\textasciitilde\textasciitilde\textasciitilde{}
谢}

衷心感谢导师xxx教授对本人的精心指导。他的言传身教将使我终生受益。感谢xxx教授，以及实验室全体老师和同窗们的热情帮助和支持！

本课题承蒙xxx基金资助，特此致谢。

最后，感谢你一直以来的努力，祝你博士毕业快乐！！！！

\clearpage

\chapterstartdual{Resume}{作者简介}

xxxx年xx月xx日出于xxxx。

xxxx年xx月考入xx大学xx院（系）xx专业，xxxx年xx月本科毕业并获得xx学学士学位。

xxxx年xx月------xxxx年xx月在xx大学xx院（系）xx学科学习并获得xx学硕士学位。

xxxx年xx月------xxxx年xx月在xx大学xx院（系）xx学科学习并获得xx学博士学位。

获奖情况：如获三好学生、优秀团干部、x奖学金等（不含科研学术获奖）。

工作经历：

\clearpage
\end{document}
